\section{Fault-to-Action Mapping: Extended Reference}
\label{sec:appendix-fault-mapping}

This section extends the round-level summary in
Table~\ref{tab:probe-outcome-mapping} by tracing each
representative fault type through the probe's observable
signature and the resulting downstream action.
Table~\ref{tab:fault-shape-mapping-full} provides a condensed
index; the paragraphs that follow supply numerical context,
evaluation cross-references, and boundary-case discussion.

\begin{table}[!htbp]
\caption{Fault-to-action index: representative fault types,
key observable ranges, and gating verdicts.}
\label{tab:fault-shape-mapping-full}
\centering
\small
\setlength{\tabcolsep}{4pt}
\begin{tabular}{@{}lcc@{}}
\toprule
Fault type & $K_{\text{eff}}/L$ & Verdict \\
\midrule
Node / shard failure     & 0.003--0.02 & \texttt{notgate} \\
One-sided backlog        & 0.003--0.08 & \texttt{notgate} \\
TTL / config drift       & 0.84--0.93  & \texttt{gate} \\
Gray failure (mixed)     & 0.015--0.43 & varies \\
Segmenter misconfiguration & ${\approx}\,1$ & \texttt{gate} \\
Zero-window segment      & --           & skip \\
\bottomrule
\end{tabular}
\end{table}

\paragraph{Node crash or shard failure}
A node or shard failure produces a small number of heavily
dirty segments (small~$B$, large~$d_{\text{big}}$),
concentrating discrepancy energy in a narrow band on the
section axis.  The concentration metric
$K_{\text{eff}}/L \approx BW/L$ falls well below the gating
threshold (measured ${\sim}0.003$--$0.02$), and the round is
admitted to triage, which outputs a compact candidate set.
When $B$ is very small, however, occupancy may drop below
$\tau_{\text{occ}}$, yielding a \texttt{unknown} verdict
instead.  This outcome is operationally benign: it signals a
concentrated pattern---precisely the regime where triage would
be most effective---rather than masking a diffuse fault.
Deployment policy can combine $\hat{D}$ and occupancy to
decide whether to attempt a budget-capped triage pass
(\S\ref{sec:eval-spatial}, \S\ref{sec:eval-ycsb} Bimodal;
\S\ref{sec:eval-dedrift}).

\paragraph{One-sided replication backlog}
When update propagation proceeds unevenly, unsynchronized keys
accumulate predominantly on one side, producing strong drift
($q$ deviating from~$0.5$) alongside a concentrated
segment-level pattern.  The de-drifted energy removes this
first-order sign bias, so both $\hat{D}$ and
$K_{\text{eff}}/L$ remain stable regardless of~$q$
(\S\ref{sec:eval-dedrift}; cross-configuration values in
Table~\ref{tab:c2-5}).  The round proceeds to sizing alert
first; the gating verdict then determines whether triage or a
fallback path follows.  Without de-drifting, the raw second
moment would produce sizing distortion on the order
of~$10^4\%$ in this scenario
(\S\ref{sec:eval-ycsb} Drift-shaped).

\paragraph{TTL drift or global configuration inconsistency}
Broad, low-rate anomalies---such as TTL mismatches or
schema-level incompatibilities---affect many segments at low
per-segment intensity, producing a diffuse energy profile.
$K_{\text{eff}}/L$ rises to ${\sim}0.84$--$0.93$, well above
the gating threshold, and the round is classified as
\texttt{gate}: triage is bypassed in favor of a conservative
fallback.  A noteworthy subtlety arises under Zipfian
segment-size skew (the Weighted Diffuse family in
\S\ref{sec:eval-ycsb}): although the record-level fault
mechanism is the same independent per-record loss, hot segments
lose disproportionately more keys, shifting the segment-level
shape toward concentration and reducing $K_{\text{eff}}/L$
to~${\sim}0.12$.  The gating metric thus tracks segment-level
shape rather than record-level fault semantics.

\paragraph{Gray failure}
Mixed fault patterns---a few heavily dirty segments embedded
in a noisy tail of mildly dirty ones---place
$K_{\text{eff}}/L$ in an intermediate range (from
${\sim}0.015$ in the base drift-stability setting to
${\sim}0.43$ under high-pressure noise coverage;
\S\ref{sec:eval-dedrift}, Table~\ref{tab:c2-5}).  Whether the
round is admitted to triage or classified as \texttt{gate}
depends on where $K_{\text{eff}}/L$ falls relative to
$\tau_{\text{gating}}$.  When noise coverage exceeds
approximately~0.3, the system approaches the gating boundary,
and roughly 4\% of conditional false alarms originate from
this interval (\S\ref{sec:appendix-far-breakdown}).  In
admitted rounds, the runtime stopping rule
(\S\ref{sec:eval-stoploss}) limits candidate expansion
(\S\ref{sec:eval-comparison}, \S\ref{sec:eval-operating}).

\paragraph{Segmenter misconfiguration}
A misconfigured segmenter effectively randomizes key-to-segment
assignment, producing near-total discrepancy
($\hat{D} \approx$ total population) with a strongly diffuse
shape ($K_{\text{eff}}/L \approx 1$).  The combination of
extreme~$\hat{D}$ and $K_{\text{eff}}/L$ near unity is rare
under normal faults, so it can serve as a dedicated alarm
condition signaling configuration error rather than data
divergence.

\paragraph{Zero-window segments}
Segments whose mapped window contains only zero-energy sections
($E_s^{(v)}{=}0$ and $\mathrm{net\_sum}_s^{(v)}{=}0$ in any
single view) can be excluded from downstream verification
without further inspection.  The derivation is given in
\S\ref{sec:appendix-sizing-zero-energy}.
